\section{The raw difference}\label{sec:raw}

How do the three populations differ before any control?

\begin{figure}
\centering
\includegraphics[width=\textwidth]{comparisons_raw.pdf}
\caption{Pooled statistics of $\Sl$, $\G$ and $\A$ with 5--95\% bands. Exponents on \PopFlightsS{},
\PopFlightsG{} and \PopFlightsA{} flights, phase statistics on the flights of
Table~\ref{tab:populations}. Values in Table~\ref{tab:raw}.}
\label{fig:raw}
\end{figure}

Group time progresses straighter over long lags than solo flights: $\Hloc(1000~\mathrm{s})$ is
\HlocThousandG{} in $\G$, \HlocThousandA{} in $\A$ and \HlocThousandS{} in $\Sl$
(Figure~\ref{fig:raw}). It also spends more time climbing, ends more searches in a climb and glides
faster and flatter, with a glide ratio of \RatioG{} against \RatioS{}. The same flights alone are no
straighter than solo flights at long lags and spend the least time in climb, but have the highest climb
rate, and sit between the two in the searches that end in a climb, the glide speed and the glide ratio. Part
of the raw difference may belong to who flies in groups, and to where and when they fly.

\begin{table}
\centering\small
\begin{tabular}{@{}l l l l@{}}
\toprule
 & solo flights $\Sl$ & group time $\G$ & same flights alone $\A$ \\
\midrule
flights (TA-MSD) & 14\,192 & 3\,342 & 2\,577 \\
$H$ (10--1000~s) & 0.946 [0.945, 0.946] & 0.943 [0.941, 0.944] & 0.946 [0.945, 0.947] \\
$H_\mathrm{loc}(1000~\mathrm{s})$ & 0.940 [0.939, 0.941] & 0.947 [0.945, 0.950] & 0.941 [0.939, 0.943] \\
time in climb (\%) & 38.7 [38.5, 38.9] & 42.1 [41.5, 42.6] & 37.6 [37.1, 38.1] \\
searches ending in a climb (\%) & 43.7 [43.4, 44.0] & 45.2 [44.5, 46.0] & 44.1 [43.3, 44.8] \\
climb rate (m/s) & 1.13 [1.12, 1.13] & 1.08 [1.06, 1.10] & 1.14 [1.12, 1.16] \\
glide speed (km/h) & 42.4 [42.3, 42.6] & 44.0 [43.2, 44.9] & 43.3 [42.8, 43.8] \\
glide ratio & 11.8 [11.7, 11.8] & 12.9 [12.6, 13.2] & 12.2 [12.0, 12.4] \\
cross-country speed (km/h) & 28.8 [28.6, 28.9] & 29.3 [28.6, 30.1] & 29.9 [29.4, 30.4] \\
\bottomrule
\end{tabular}

\caption{Pooled statistics of the populations of Table~\ref{tab:populations}, with 5--95\% bands. The
first row is the number of flights of the exponents; the other rows rest on the flights with phases
(Table~\ref{tab:populations}).}
\label{tab:raw}
\end{table}

\subsection{Longer lags and the size of the group}

How do the curves go over lags of hours, and with the number of pilots of the group?

A piece of at least 1250~s gives the lags up to 1000~s. A longer lag $\tau$ needs pieces of at least
$\tau/0.8$~s, so a curve to a long lag is measured on the long pieces of its population only, and the
population does not change along it. The populations of one panel of Figure~\ref{fig:long} share the
longest lag $T$, the longest at which every one of them keeps at least \LongMinFlights{} flights, and are
measured on the pieces of at least $T/0.8$~s (dashed); the thin lines are the curves of all the pieces
up to 1000~s. The last panel gives the flights left at each lag, and Table~\ref{tab:long} the flights of
each curve. $\Hloc$ is read from $\MSD$ over the lags within $\pm0.25$ decades.

\begin{figure}
\centering
\includegraphics[width=\textwidth]{longlags.pdf}
\caption{$\Hloc(\tau)$ with an inset of $\MSD(\tau)$ in $\mathrm{m}^2$. Top left: $\Sl$, $\G$ and $\A$, lags up to
\LongTauPooled~s. Top right: $\Sl$ and group time by the number of pilots of the group, up to
\LongTauSize~s. Bottom left: $\Sl$ and the same flights alone, by the size of the group the flight
had, up to \LongTauAlone~s. Each population has two curves of its colour. Thin: all its pieces, up to
1000~s, with its flights in the legend. Dashed: only its pieces long enough for the longest lag of the
panel, a subset of the thin ones, measured over the same population at every lag; 5--95\% bands.
Bottom right: flights with pieces of at least $\tau/0.8$~s, one line per population (dashed for the
flights alone); dotted, the minimum of \LongMinFlights{} flights.}
\label{fig:long}
\end{figure}

\begin{table}
\centering\small
\begin{tabular}{@{}l r r@{}}
\toprule
population & flights up to 1000~s & flights of the long series \\
\midrule
\multicolumn{3}{@{}l}{\emph{pooled populations, lags up to 9\,350~s}} \\
solo flights $\Sl$ & 14\,192 & 2\,436 \\
group time $\G$ & 3\,342 & 130 \\
same flights alone $\A$ & 2\,577 & 353 \\
\midrule
\multicolumn{3}{@{}l}{\emph{group time by the size of the group, lags up to 1\,230~s}} \\
solo flights $\Sl$ & 14\,192 & 13\,988 \\
$\G$, 2 pilots & 2\,577 & 2\,494 \\
$\G$, 3 pilots & 1\,022 & 950 \\
$\G$, 4 pilots & 398 & 348 \\
$\G$, 5--6 pilots & 305 & 258 \\
$\G$, 7+ pilots & 144 & 107 \\
\midrule
\multicolumn{3}{@{}l}{\emph{the same flights alone, by size, lags up to 4\,150~s}} \\
solo flights $\Sl$ & 14\,192 & 7\,916 \\
$\A$, flights of 2 pilots & 2\,001 & 945 \\
$\A$, flights of 3 pilots & 762 & 356 \\
$\A$, flights of 4 pilots & 293 & 129 \\
$\A$, flights of 5--6 pilots & 225 & 102 \\
\bottomrule
\end{tabular}

\caption{Flights of each curve of Figure~\ref{fig:long}: those of all the pieces, up to 1000~s, and those
of the long series, with pieces of at least $T/0.8$~s for the longest lag $T$ of the panel.}
\label{tab:long}
\end{table}

\paragraph{Pooled.} On the long pieces group time stays straight longer. At \LongLagMid~s
$\Hloc$ is \LongPooledGMid{} \LongPooledGMidBand{} in $\G$, against \LongPooledSMid{} in $\Sl$ and
\LongPooledAMid{} in $\A$; at \LongLagEnd~s it is \LongPooledGEnd{} \LongPooledGEndBand{} against
\LongPooledSEnd{} and \LongPooledAEnd{}. For every population $\Hloc$ falls with the lag as the motion
saturates, and $\A$ is below $\Sl$. The long series of $\G$ rests on few flights
(Table~\ref{tab:long}), and its band is wide.

\paragraph{By size.} Group time is straighter than solo flights at every size, though the band of 7+ pilots reaches $\Sl$: at \LongLagSize~s $\Hloc$
is \LongSizeS{} in $\Sl$ and \LongSizeTwo{} \LongSizeTwoBand{} for 2 pilots, \LongSizeThree{},
\LongSizeFour{}, \LongSizeFiveSix{} and \LongSizeSevenPlus{} \LongSizeSevenPlusBand{} for 3, 4, 5--6 and 7+
pilots. Two pilots are the least straight; from three on the sizes are within each other's bands.

\paragraph{Alone, by size.} The solo time of the flights of groups is not straighter than solo flights:
at \LongLagAlone~s \LongAloneTwo{} \LongAloneTwoBand{} for flights that had groups of two, against
\LongAloneS{} in $\Sl$; \LongAloneFiveSix{} \LongAloneFiveSixBand{} for 5--6 pilots, whose band reaches $\Sl$,
and \LongAloneSevenPlus{} \LongAloneSevenPlusBand{} for 7+.

\paragraph{What we keep.} A difference to explain, largest at long lags and in the glides: group time
against solo flights, while the same flights alone are no straighter than solo flights; on the long pieces group time keeps its
progress to lags of hours, at every size of the group. The next sections ask where the difference
comes from.
